\documentclass[11pt]{article}
\usepackage[margin=1in]{geometry}
\usepackage{amsmath,amssymb,amsthm,booktabs}
\newtheorem{theorem}{Theorem}
\newtheorem{proposition}{Proposition}
\begin{document}
\title{Further discontinuities from four-point ratio-graph certificates}
\author{}\date{}\maketitle

\paragraph{Endpoints.}
We consider the next intermediate threshold after $\alpha_2$,
\[
\alpha_3^5-\alpha_3^4-\alpha_3^2-\alpha_3-2=0,
\qquad \alpha_3=1.739234341260724\ldots,
\]
and further members of a branching family beginning at $\sqrt[3]4$,
\[
\tau_1^4-\tau_1^2-4=0,\qquad
\tau_1=\sqrt{\frac{1+\sqrt{17}}2}=1.600485180440241\ldots.
\]
The next two are defined by
\[
\tau_2^5=\tau_2^3+\tau_2^2+4,\qquad
\tau_3^6=\tau_3^4+\tau_3^3+\tau_3^2+4.
\]
Each is the unique positive root of its polynomial. Exact fractions and
the final jump comparisons are recorded in \texttt{family\_discontinuities.json}.

\begin{theorem}
The exact upper certificates give
\[
S_2(\alpha_3)\le\frac{3050935043544191377}{6000000000000000000},
\qquad
S_2(\tau_1)\le\frac{347114271361950569}{750000000000000000}.
\]
Together with the uniform right-side constructions below, these prove
right-hand discontinuities at these exact algebraic endpoints:
\[
S_2(\alpha_3+)-S_2(\alpha_3)>0.0005649589651,
\qquad
S_2(\tau_1+)-S_2(\tau_1)>0.0004019739059.
\]
The same method also proves discontinuities at $\tau_2$ and $\tau_3$:
\[
S_2(\tau_2+)-S_2(\tau_2)>0.0002286755200,\qquad
S_2(\tau_3+)-S_2(\tau_3)>0.0001522007113.
\]
Their exact fractions and constructions are given below.
\end{theorem}
The following decimals round upper bounds upward and lower bounds downward:
\begin{center}
\begin{tabular}{lrrr}
\toprule
endpoint&upper at endpoint&uniform right lower&jump lower\\\midrule
$\alpha_3$&0.508489173925&0.509054132889&0.0005649589651\\
$\tau_1$&0.462819028483&0.463221002388&0.0004019739059\\
$\tau_2$&0.464038146566&0.464266822085&0.0002286755200\\
$\tau_3$&0.464794534376&0.464946735086&0.0001522007113\\
\bottomrule
\end{tabular}
\end{center}

\paragraph{Upper certificates.}
All certificates use four iid points, pairwise ratio colors, and the
additive triple labels recording the winning event exactly. For $x<y$,
the color of the pair is the interval containing $y/x$. At $c=\alpha_3$
the six cutoffs are
\[
\frac{c^2+c+2}{c(c+2)},\quad
\frac{c+2}{2c},\quad \frac2c,\quad
\frac{c+2}{c^2},\quad\frac1{c-1},\quad c.
\]
They are ratios suggested by the intermediate ladder. At $c=\tau_1$
the cutoffs are
\[
\frac{c+2}{2c},\quad\frac2c,\quad\frac{c+2}{c^2},\quad
\frac{c^2+c+2}{c^3},\quad\frac4{c^2},\quad c.
\]
These combine the original ladder thresholds with the branching threshold
$4/c^2$. The listed cutoffs are in increasing order at their endpoints.

The flag families are $(t,f)=(0,2),(1,2),(2,3)$. There are seven pair
colors, so the last family has seven root types. The exact cover sizes
and PSD block dimensions are
\[
\begin{array}{c|r|r|l}
&\text{patterns}&\text{excluded prefixes}&\text{block dimensions}\\\hline
\alpha_3&1995&18611&7,14,47,55,42,40,40,32,32\\
\tau_1&1974&18336&7,14,72,33,30,33,38,35,47.
\end{array}
\]

For every excluded prefix, the cover stores a nonnegative nonzero
combination of the prescribed strict linear inequalities whose coefficient
vector in positive gap coordinates is nonnegative. All these contradictions
are checked in the exact algebraic field of the endpoint. Field elements
are coefficient vectors modulo the monic endpoint polynomial; signs use
rational isolation of its unique positive root. The full binary partition
of retained patterns and excluded prefixes is checked independently.

For each retained pattern $P$, the verifier evaluates
\[
\operatorname{obj}(P)+\sum_\sigma\langle Q_\sigma,M_\sigma(P)\rangle,
\qquad Q_\sigma=L_\sigma L_\sigma^T/10^{18},
\]
with integer $L_\sigma$. Its exact maximum is the upper bound. Conditional
independence makes each flag term nonnegative in expectation; thus the
bound holds for all iid laws. Atomless approximation transfers the result
to arbitrary nonnegative laws, as in the earlier cube-root certificate.
The proof does not rely on numerical feasibility or solver convergence.

\paragraph{Uniform lower family at $\alpha_3$.}
We use the existing repository construction with masses on
$0,a_0,a_1,a_2,a_3,1$. For $\alpha_3<c<2$ set
\[
r_0=u_0=1/c,\quad r_{j+1}=(r_0+r_j)/c,\quad
u_{j+1}=(u_0+u_j+1)/c,\quad
\delta=\frac{c-1-r_3}{2u_3},\quad a_j=r_j+\delta u_j.
\]
Then $\delta>0$, $ca_0=1+\delta$,
$ca_{j+1}=a_0+a_j+\delta$, and
$a_3=(c-1+r_3)/2<c-1$. These imply all selected successes uniformly.
The six exact masses are
\[
\frac{(309618399286,175410569272,48113766390,
48113766390,48113766390,370629732273)}{1000000000001}.
\]
Count the following disjoint classes of ordered successes: positive
target with at least one zero summand; target $1$ with two positive
summands, at least one smaller than $1$; and an intermediate target
$a_k$ with summands $a_0,a_j$, $j<k$. Recursive zero replacement gives
the exact uniform bound
\[
L_\alpha=
\frac{493944836144780432326982586906209328}
{970318880118646973873431854824794345}
=0.5090541328891619\ldots\qquad(c>\alpha_3).
\]
The masses and selected-event polynomial come from the repository's
\texttt{multinacci\_families.json}; the new contribution is the separating
endpoint upper certificate.

\paragraph{Uniform lower family at $\tau_1$.}
For $\tau_1<c<\varphi$ start from the limiting locations
\[
1,\quad \frac2c,\quad \frac{c+2}{c^2},\quad
\frac4{c^2},\quad \frac{c^2+4}{c^3}.
\]
They are strictly ordered and the last is below $c$. To make all
needed comparisons strict, take sufficiently small $\varepsilon>0$ and set
\[
x_0=1,\quad x_1=2/c+\varepsilon,\quad
x_2=(1+x_1)/c+\varepsilon,\quad
x_3=2x_1/c+\varepsilon,\quad
x_4=(1+x_3)/c+\varepsilon.
\]
The ordering and $x_4<c$ persist by continuity. The only required
comparison not immediate from the recursion, $1+x_2<cx_3$, has limiting
slack $(2-c)(c-1)/c^2>0$. Thus $x_0+x_i<cx_j$ for $i<j$ and also
$2x_1<cx_j$ for $j=3,4$.

Take positive masses
\[
(p_0,p_1,p_2,p_3,p_4)
=(515633,163083,89902,115691,115691)/10^6.
\]
The selected within-scale probability is
\[
B=p_0^2\sum_{j=1}^4p_j
+2p_0\sum_{1\le i<j\le4}p_i p_j+p_1^2(p_3+p_4)
=\frac{224225230418916583}{10^{18}}.
\]
Recursive scale replacement with $w=560802159/10^9$ gives
\[
L_\tau=\frac{1-(1-B)w^2}{3-3w+w^2}
=\frac{756019563001067760932969902864122823}
{1632092584539061281000000000000000000}
=0.463221002388528\ldots.
\]
Monotonicity extends the bound from $(\tau_1,\varphi)$ to every
$c>\tau_1$.

\begin{proposition}[The branching family]
Define $b_0=4/c^2$ and $b_{j+1}=(1+b_j)/c$. The condition $b_j<c$
opens at the unique positive root $\tau_j$ of
\[
\boxed{\tau_j^{j+3}=4+\sum_{k=2}^{j+1}\tau_j^k},\qquad j\ge0.
\]
The empty sum gives $\tau_0=\sqrt[3]4$. If $c_n$ are the original
ladder thresholds, then $c_{j+3}<\tau_j<c_{j+4}$ and
$\tau_j\uparrow\varphi$.
\end{proposition}
\begin{proof}
Induction gives
$b_j=(4+\sum_{k=2}^{j+1}c^k)/c^{j+2}$. For the original ladder
$r_0=1$, $r_{k+1}=(1+r_k)/c$, and $1<c<\varphi$,
\[
b_0-r_3=\frac{(2-c)(c-1)}{c^3}>0,\qquad
r_4-b_0=\frac{(c-2)(c^2-c-1)}{c^4}>0.
\]
Iterating the increasing affine map gives $r_{j+3}<b_j<r_{j+4}$.
Evaluating at the respective ladder thresholds proves the interlacing;
their convergence to $\varphi$ gives the limit.
\end{proof}

For the later members we use the sorted ratio cutoffs
\[
\left\{\frac{c+2}{2c},r_1,\ldots,r_{j+2},b_0,\ldots,b_{j-1},c\right\}.
\]
There are eight cutoffs at $\tau_2$ and ten at $\tau_3$. The additive
triple signs and the four-point flag families remain the same.

More generally the template $1,r_1,r_2,b_0,\ldots,b_j$, slightly perturbed,
gives a uniform lower bound for every $c>\tau_j$. For masses
$p_0,\ldots,p_{j+3}$ its guaranteed one-scale probability is
\[
p_0^2\sum_{i=1}^{j+3}p_i
+2p_0\sum_{1\le i<k\le j+3}p_i p_k
+p_1^2\sum_{k=3}^{j+3}p_k.
\]
To justify strict feasibility for all these templates, use the same
$x_0,x_1,x_2,x_3$ as above and continue
$x_{k+1}=(1+x_k)/c+\varepsilon$ for $k\ge3$.
For $\tau_j<c<\varphi$, sufficiently small $\varepsilon>0$ preserves
ordering and $x_{j+3}<c$. All the selected successes are then strict.

For completeness, put mass $w$ on this block and recursively put mass
$1-w$ on a sufficiently smaller copy. The selected contributions are
$w^3B$, $2w^2(1-w)$, $w(1-w)^2$, and the recursive contribution
$(1-w)^3 S$. Finite truncations therefore approach
\[
L(B,w)=\frac{1-(1-B)w^2}{3-3w+w^2}.
\]
The scale separation can be chosen uniformly across levels for each fixed
$c>\tau_j$, since the positive block has maximum less than $c$ times its
minimum. Monotonicity covers $c\ge\varphi$.
The exact rational masses and $w$ for the later endpoints are saved in
\texttt{branch2\_lower.json} and \texttt{branch3\_lower.json}.

\paragraph{The next branching discontinuity.}
At $\tau_2=1.6076824592666596\ldots$, defined by
$\tau_2^5=\tau_2^3+\tau_2^2+4$, the exact upper certificate gives
\[
S_2(\tau_2)\le\frac{928076293130715131}{2000000000000000000}
<0.464038146566.
\]
The uniform right lower is
\[
L_2=\frac{18872601477188494350408465591458906090728633}
{40650334203112904100000000000000000000000000}
>0.464266822085.
\]
Consequently $S_2(\tau_2+)-S_2(\tau_2)>0.0002286755200$.
The cover contains 4162 patterns and 47968 excluded prefixes, with PSD
block dimensions $9,18,112,39,40,38,38,43,52,45,73$.

\paragraph{Discontinuity at $\tau_3$.}
At $\tau_3=1.6118224692844974\ldots$, the original ten-cutoff
coloring suffices. The saved cover has 7542 patterns and 103331 exact
excluded prefixes, with PSD block dimensions
\[
11,22,160,47,46,46,48,46,46,55,64,57,105.
\]
A longer optimization of this same SDP gives the exact certificate
\[
S_2(\tau_3)\le
\frac{929589068751260629}{2000000000000000000}
<0.464794534376.
\]
The fixed masses in \texttt{branch3\_lower.json}, using the uniform
branch construction proved above, give
\[
L_3=\frac{18854601744361333727955719879293630392307573}
{40552175811785600100000000000000000000000000}
>0.464946735086.
\]
Their exact positive difference is
\[
L_3-U_3=
\frac{6172070003747467725548707159398942307573}
{40552175811785600100000000000000000000000000}
>0.0001522007113.
\]
Allowing the full solve more time produced a separating certificate
with the same sampled points, colors, and flag families.
The search used a 600-second solver budget (about 644 seconds elapsed in
the recorded run). The final proof uses only the rational Gram factors
and the exact cover, independently of the solver status or tolerances.

As a separate experiment, adding the branch-atom ratio
$(c^2+4)/(4c)$ as an eleventh cutoff and optimizing Gram factors directly
also gives a separating exact upper:
\[
S_2(\tau_3)\le\frac{557935142150665721}{1200000000000000000}.
\]
This weaker alternative uses 9625 patterns and 151158 excluded prefixes.
Its cover reuses old exclusions by appending the new comparison forms;
the verifier checks that the stored order is a permutation of all the
required forms. The stronger theorem above uses the original coloring.

\paragraph{Reproduction.}
The alternative eleventh-cutoff cover is stored with Git LFS:
\begin{verbatim}
branch3_extra_branch_ratio_cover.json
\end{verbatim}
Before using that alternative cover, run from the repository root:
\begin{verbatim}
git lfs install
git lfs pull
\end{verbatim}
The original ten-cutoff $\tau_3$ cover is stored in regular Git.

From the repository root, the standard-library command
\begin{verbatim}
python3 graph_experiment/verify_family_discontinuities.py
\end{verbatim}
checks all saved full upper certificates, recomputes the exact lower
probabilities from the saved masses, and verifies the positive gaps.
For numerical regeneration, use the pinned search environment and run
\begin{verbatim}
python graph_experiment/family_endpoint_exact.py alpha3
python graph_experiment/family_endpoint_exact.py branch1
python graph_experiment/family_endpoint_exact.py branch2
python graph_experiment/family_endpoint_exact.py branch3 --seconds 600
\end{verbatim}
The exact search reuses saved covers; in a fresh directory without JSON
files it regenerates every algebraic exclusion. Numerical four-color
tests did not separate either endpoint. Adding the two further
ladder/branch cutoffs gave the successful six-color-cutoff formulations.
\end{document}
